\section{Results}
\label{sec:results}

\subsection{Double-differential cross section}
\label{sec:results-xsec}
Figure~\ref{fig:xsec} shows the unfolded $d^{2}\sigma/(d\pT\,d\pPar)$ and its
1D projections. The total cross section is
\SI{\sigTwoD}{cm^2/nucleon} (published: \SI{\sigTwoDpaper}{cm^2/nucleon}).

\begin{figure}[H]
  \centering
  \includegraphics[width=0.62\textwidth]{MEFHC_5iter_fig13.png}\\
  \includegraphics[width=0.49\textwidth]{MEFHC_5iter_xsec_proj_pt.png}
  \includegraphics[width=0.49\textwidth]{MEFHC_5iter_xsec_proj_pz.png}
  \caption{Unfolded 2D cross section (top, Fig.~13 style) and its $\pT$/$\pPar$
  projections (bottom).}
  \label{fig:xsec}
\end{figure}

\subsection{Comparison to the published measurement}
\label{sec:results-compare}
\begin{table}[h]
  \centering
  \begin{tabular}{@{}lc@{}}
    \toprule
    Quantity (205 reported bins) & Value \\
    \midrule
    $\sigma_{\mathrm{tot}}(\mathrm{ours})/\sigma_{\mathrm{tot}}(\mathrm{paper})$ & $\ratioTot$ \\
    Median bin ratio (ours/paper) & $\medianBinRatio$ \\
    Bins within \SI{5}{\percent}/\SI{10}{\percent}/\SI{20}{\percent} & $\binsFive/\binsTen/\binsTwenty\,\%$ \\
    $\chi^{2}/\mathrm{ndf}$, paper covariance & $\chiPaper$ \\
    $\chi^{2}/\mathrm{ndf}$, combined covariance (paper $+$ ours) & $\chiCombined$ \\
    Pull mean / RMS (paper covariance) & $\pullMean/\pullRMS$ \\
    \bottomrule
  \end{tabular}
  \caption{Agreement with arXiv:2106.16210. Per-bin pulls
  (Fig.~\ref{fig:pull}) are sub-$1\sigma$.}
  \label{tab:compare}
\end{table}

\begin{figure}[H]
  \centering
  \includegraphics[width=0.7\textwidth]{MEFHC_5iter_pull_full.png}
  \caption{Per-bin pull $(\mathrm{ours}-\mathrm{paper})/\sigma$ on the
  $14\times16$ grid and its distribution.}
  \label{fig:pull}
\end{figure}

\begin{figure}[H]
  \centering
  \includegraphics[width=0.99\textwidth]{MEFHC_5iter_xsec_paper_pt_slices.png}
  \caption{Bin-by-bin comparison to the published measurement: the unfolded
  cross section (points) overlaid on the published data (histogram) in slices
  of $\pT$ versus $\pPar$.}
  \label{fig:ptslices}
\end{figure}

\begin{figure}[H]
  \centering
  \includegraphics[width=0.99\textwidth]{MEFHC_5iter_xsec_paper_pz_slices.png}
  \caption{Same comparison as Fig.~\ref{fig:ptslices}, sliced in $\pPar$
  versus $\pT$.}
  \label{fig:pzslices}
\end{figure}

\subsection{Anatomy of the $\chi^{2}=3.66$ tension}
\label{sec:results-tension}
The central values agree (pull RMS $\pullRMS$, so a diagonal-only $\chi^{2}/\mathrm{ndf}$
would be $0.37$), yet the full paper-covariance $\chi^{2}/\mathrm{ndf}$ is $\chiPaper$:
a $10\times$ inflation living entirely in the off-diagonal structure of the
published covariance. The eigenmode decomposition (Fig.~\ref{fig:tension}) shows
the tension is carried by moderate-$\lambda$ modes with $5$--$6\sigma$ eigen-pulls
(the 10 smallest-$\lambda$ modes carry only \SI{3}{\percent}; keeping the 73
best-measured directions still gives $1.42$), so it is \emph{not} a small-$\lambda$
inversion artifact. It is \emph{shape}, not normalisation
($\chi^{2}_{\mathrm{norm}}=0.10$ vs $\chi^{2}_{\mathrm{shape}}=750$), localised to
the low-$\pT$ cross-section peak ridge where the flux and muon-energy bands
dominate. A $\sim1$-unit methodological band comes from the GBDT estimator
(exact-GBT $\chiPaper$ vs HistGBT $2.70\pm0.04$ vs LightGBM $2.65$; iteration count is
negligible, $5\to10$ iters move it by $+0.01$). Full treatment:
App.~\ref{sec:tension}. The eigenmode-localisation maps and the additional
model-comparison ratio maps are shown in
Figs.~\ref{fig:tension}--\ref{fig:tensionmodes}.
% NOTE: the estimator-band value "exact-GBT 3.66" below is the same number as
% the paper-covariance chi2 (\chiPaper); they coincide by construction.

\begin{figure}[H]
  \centering
  \includegraphics[width=0.49\textwidth]{tension_spectrum.png}
  \includegraphics[width=0.49\textwidth]{tension_chi2_map.png}
  \caption{Tension anatomy: eigenvalue spectrum with cumulative $\chi^{2}$ (left)
  and the per-bin $\chi^{2}$-contribution map (right).}
  \label{fig:tension}
\end{figure}

\begin{figure}[H]
  \centering
  \includegraphics[width=0.9\textwidth]{tension_mode_maps.png}
  \caption{Bin-space maps of the leading tension-carrying eigenmodes of the
  published covariance: the $5$--$6\sigma$ eigen-pulls live on the low-$\pT$
  cross-section-peak ridge where the flux and muon-energy bands dominate.}
  \label{fig:tensionmodes}
\end{figure}

\subsection{Comparison to GENIE MINERvA Tune v1}
\label{sec:results-model}
Both the published data and our unfolded result are compared to the GENIE
2.12.6~\cite{Andreopoulos:2009rq} MINERvA Tune v1 prediction shipped with the
ancillary (Fig.~\ref{fig:model}). The tune
under-predicts the total by $\sim\SI{9}{\percent}$
($\sigma_{\mathrm{tune}}/\sigma_{\mathrm{data}}=0.91$), concentrated in the
low-$\pT$ peak, and is strongly disfavoured by the data covariance
($\chi^{2}/\mathrm{ndf}$: data vs tune $33.0$, ours vs tune $26.5$). Our unfolded
result tracks the data far more closely than the model does.
(\texttt{compare\_to\_models.py}.) The published analysis~\cite{MINERvA:2021owq}
additionally compares to NuWro~\cite{Golan:2012wx}, GiBUU~\cite{Buss:2011mx} and
NEUT~\cite{Hayato:2009zz}; because OmniFold delivers a per-event reweighted
truth ensemble, any such generator prediction can be compared to our result in
\emph{any} projection or dimension without re-unfolding --- a capability we
exploit for the 3D extension (\S\ref{sec:3d}).

\begin{figure}[H]
  \centering
  \includegraphics[width=0.85\textwidth]{model_comp_projections.png}
  \caption{Our result (dashed) on top of the published data (black); MINERvA
  Tune v1 (orange) dips below in the peak.}
  \label{fig:model}
\end{figure}

\begin{figure}[H]
  \centering
  \includegraphics[width=0.9\textwidth]{model_comp_ratio_maps.png}
  \caption{Bin-by-bin ratio maps on the $(\pT,\pPar)$ grid: ours/paper (left)
  and Tune~v1/paper (right). The tune's deficit is concentrated in the
  low-$\pT$ peak; our result is flat near unity.}
  \label{fig:modelmaps}
\end{figure}
